\clearpage\pagestyle{otMain}\setlength{\otBlocEspacement}{2.5mm}\fontsize{11}{15}\selectfont\corrigesChapitre{17}{Espace, pyramides et cônes}
\begin{blocLivre}{cor-17-1}\begin{corrige}{1}{Nommer les éléments}\demarche{Distinguer base et faces latérales.}
Le sommet est $\res S$, la base est le carré $\res{ABCD}$. Il y a $\res4$ faces latérales triangulaires : SAB, SBC, SCD et SDA. Avec la base, le solide a cinq faces.
\label{fig-cor-17-1}\figureGeo[corrige]{sol-01}
\end{corrige}\end{blocLivre}
\begin{blocLivre}{cor-17-2}\begin{corrige}{2}{Une pyramide à base carrée}\demarche{Calculer d’abord l’aire de base.}
L’aire de la base vaut $6^2=36$ cm$^2$. Le volume est $V=36\times8/3=\res{96\ \mathrm{cm}^3}$. La hauteur donnée est perpendiculaire à la base.
\label{fig-cor-17-2}\figureGeo[corrige]{sol-02}
\end{corrige}\end{blocLivre}
\begin{blocLivre}{cor-17-3}\begin{corrige}{3}{Un cône}\demarche{Appliquer la formule avec le rayon.}
$V=\pi\times3^2\times5/3=\res{15\pi\ \mathrm{cm}^3}$. Numériquement, $V\approx\res{47{,}1\ \mathrm{cm}^3}$. Garder $\pi$ pour la valeur exacte.
\label{fig-cor-17-3}\figureGeo[corrige]{sol-03}
\end{corrige}\end{blocLivre}
\begin{blocLivre}{cor-17-4}\begin{corrige}{4}{Du diamètre au rayon}\demarche{Diviser le diamètre par deux.}
Le rayon vaut $8/2=4$ cm. Donc $V=\pi\times4^2\times9/3=\res{48\pi\ \mathrm{cm}^3}$. Utiliser 8 comme rayon multiplierait l’aire de base par 4.
\label{fig-cor-17-4}\figureGeo[corrige]{sol-04}
\end{corrige}\end{blocLivre}
\begin{blocLivre}{cor-17-5}\begin{corrige}{5}{Lire trois coordonnées}\demarche{Suivre les axes dans l’ordre abscisse, ordonnée, altitude.}
Le sommet opposé utilise les trois dimensions : $\res{G(6;4;3)}$. H a la même abscisse 0 et la même ordonnée 4 que son pied, avec altitude 3 : $\res{H(0;4;3)}$.
\label{fig-cor-17-5}\figureGeo[corrige]{sol-05}
\end{corrige}\end{blocLivre}
\begin{blocLivre}{cor-17-6}\begin{corrige}{6}{Placer un point}\demarche{Reporter les trois coordonnées dans leur direction.}
Depuis O, avancer de 2 unités selon $x$, puis de 3 selon $y$, puis monter d’une unité selon $z$. Les pointillés de construction se rencontrent en $\res{M(2;3;1)}$, à l’intérieur du pavé.
\label{fig-cor-17-6}\figureGeo[corrige]{sol-06}
\end{corrige}\end{blocLivre}
\begin{blocLivre}{cor-17-7}\begin{corrige}{7}{Base triangulaire}\demarche{Distinguer hauteur de base et hauteur du solide.}
L’aire de base est $8\times3/2=12$ cm$^2$. La hauteur du solide vaut 10 cm, différente de la hauteur du triangle de base. $V=12\times10/3=\res{40\ \mathrm{cm}^3}$.
\label{fig-cor-17-7}\figureGeo[corrige]{sol-07}
\end{corrige}\end{blocLivre}
\begin{blocLivre}{cor-17-8}\begin{corrige}{8}{Retrouver une hauteur}\demarche{Isoler la hauteur dans la formule.}
$150=45h/3=15h$. Donc $h=150/15=\res{10\ \mathrm{cm}}$. La formule donne une longueur car on divise un volume par une aire.
\label{fig-cor-17-8}\figureGeo[corrige]{sol-08}
\end{corrige}\end{blocLivre}
\begin{blocLivre}{cor-17-9}\begin{corrige}{9}{Un patron de pyramide}\demarche{Faire correspondre les côtés qui seront assemblés.}
Tracer un carré de côté 4 cm. À l’extérieur, construire sur chacun de ses côtés un triangle isocèle de côtés 5, 5 et 4 cm. Les quatre bases des triangles coïncident avec les côtés du carré ; leurs sommets se rejoindront au pliage. Le patron comporte un carré et quatre triangles.
\label{fig-cor-17-9}\figureGeo[corrige]{sol-09}
\end{corrige}\end{blocLivre}
\begin{blocLivre}{cor-17-10}\begin{corrige}{10}{Un patron de cône}\demarche{Relier la perspective au patron par les longueurs.}
Pour la perspective, tracer une ellipse pour la base, placer le sommet au-dessus du centre et le relier au bord ; représenter en pointillés les parties cachées. Le bord courbe du secteur doit mesurer le périmètre de base $2\pi\times3=6\pi$ cm. Le cercle de rayon 5 a pour périmètre $10\pi$ ; il faut donc $6/10$ de tour, soit $\res{216^\circ}$. Ajouter le disque de rayon 3 cm. La génératrice, et non la hauteur, est le rayon du secteur.
\label{fig-cor-17-10}\figureGeo[corrige]{sol-10}
\end{corrige}\end{blocLivre}
\begin{blocLivre}{cor-17-11}\begin{corrige}{11}{Une génératrice calculée}\demarche{Séparer le triangle de section et la formule du volume.}
La hauteur, le rayon et la génératrice forment un triangle rectangle. La génératrice vaut $\sqrt{12^2+5^2}=\res{13}$ cm. Le volume est $\pi\times5^2\times12/3=\res{100\pi\ \mathrm{cm}^3}$. La génératrice ne remplace pas la hauteur dans cette formule.
\label{fig-cor-17-11}\figureGeo[corrige]{sol-11}
\end{corrige}\end{blocLivre}
\begin{blocLivre}{cor-17-12}\begin{corrige}{12}{Comparer un cylindre et un cône}\demarche{Utiliser la base et la hauteur communes.}
Le cylindre a pour volume $\pi\times16\times6=\res{96\pi}$ cm$^3$. Le cône a le tiers, soit $\res{32\pi}$ cm$^3$. Leur différence vaut $\res{64\pi\ \mathrm{cm}^3}$.
\label{fig-cor-17-12}\figureGeo[corrige]{sol-12}
\end{corrige}\end{blocLivre}
\begin{blocLivre}{cor-17-13}\begin{corrige}{13}{Deux unités}\demarche{Choisir les décimètres pour obtenir directement des litres.}
On choisit les dm : $r=4$ dm et $h=12$ dm. $V=\pi\times4^2\times12/3=64\pi$ dm$^3$, donc $64\pi$ L, soit environ $\res{201{,}1\ \mathrm L}$. Les unités sont unifiées avant le calcul.
\label{fig-cor-17-13}\figureGeo[corrige]{sol-13}
\end{corrige}\end{blocLivre}
\begin{blocLivre}{cor-17-14}\begin{corrige}{14}{Un assemblage}\demarche{Additionner les volumes des deux parties.}
Le pavé a pour volume $8\times6\times4=192$ cm$^3$. La pyramide a pour volume $(8\times6)\times3/3=48$ cm$^3$. Les deux intérieurs ne se chevauchent pas : le total vaut $\res{240\ \mathrm{cm}^3}$.
\label{fig-cor-17-14}\figureGeo[corrige]{sol-14}
\end{corrige}\end{blocLivre}
\begin{blocLivre}{cor-17-15}\begin{corrige}{15}{Un remplissage partiel}\demarche{Appliquer la fraction au volume, sans la transférer à la hauteur.}
La capacité vaut $\pi\times36\times10/3=120\pi$ cm$^3$. L’eau occupe $3/4$ de cette capacité, soit $\res{90\pi\ \mathrm{cm}^3}\approx\res{282{,}7\ \mathrm{mL}}$. On ne peut pas conclure à une hauteur de 7,5 cm : la section du cône change avec la hauteur, donc volume et hauteur ne sont pas proportionnels.
\label{fig-cor-17-15}\figureGeo[corrige]{sol-15}
\end{corrige}\end{blocLivre}
